%============================================================
% CH2 — SOLUTION SKETCHES FOR THE PRACTICE SET   (tag g2s)
% Graph model, algorithm and complexity for each problem.
% Sketches, not full code: the point is to check your MODEL.
% Every sketch was implemented in C++, reproduces the official samples
% (LeetCode: the examples) and was compared with a brute force or a
% checker on random small inputs. Order follows problem_set.tex.
%============================================================
\subsection{Solution Sketches}
\label{ssec:ch2-sketches}

\begin{warningbox}[Read a sketch only after a real attempt]
Each sketch gives the \emph{design} --- the graph model, what the DFS computes on the way in and on the way out, the complexity and the one trap of the problem --- in a few lines, so you can compare it with your own model. Writing and debugging the code is left to you. If your model differs but you can argue why it answers the question, it is probably also right. ``Brute force'' below means an independent slow program (exhaustive search, simulation or a closed formula) that the fast solution was compared with on random small inputs.
\end{warningbox}

\begingroup
\providecommand{\sk}{}
\renewcommand{\sk}[2]{\par\smallskip\noindent\textbf{\textcolor{deepblue}{#1}}\quad #2\par}

\begin{solutionbox}[Tier 1 --- Warm-up]
\sk{P2.01 Gold Rush (CF 1829D).}{Implicit tree: the children of a pile $x$ with $3 \mid x$ are $x/3$ and $2x/3$; a pile never grows. \texttt{can(x)}: true if $x = m$, false if $x < m$ or $3 \nmid x$, else \texttt{can(x/3) || can(2x/3)}. Every call divides by $3$, and $3^{15} > 10^7$, so the depth is at most $14$ and there are fewer than $2^{15}$ calls: plain recursion is fine. Samples: \texttt{YES YES NO NO YES YES NO YES YES NO NO}; checked against the closed form ``$n = m \cdot 3^a / 2^b$ with $b \le a$'' on $10\,000$ random cases.}
\sk{P2.02 Huge Pile (CF 2184C).}{The split tree of $n$ has up to $n$ leaves, but by induction every pile at depth $d$ is $\lfloor n/2^d \rfloor$ or $\lceil n/2^d \rceil$. Keep the set of distinct pile sizes per level (at most $2$), go down until $k$ appears (answer $d$) or all piles are $1$ (answer $-1$). \bigO{\log n} per test. Samples: $2, 1, -1, 29$; checked against a DFS over the full, unmerged split tree on $8000$ random cases.}
\sk{P2.03 Permutation Transformation (CF 1490D).}{The tree is the recursion itself: \texttt{build(l, r, dep)} finds the maximum position $m$ of $a[l..r]$, records $d_m = \textit{dep}$ and recurses on $[l, m-1]$ and $[m+1, r]$ with $\textit{dep}+1$ --- the depth is pushed \emph{down} as a parameter, no tree is stored. \bigO{n^2} worst case, fine for $n \le 100$. Samples: \texttt{1 0 2 3 1}, \texttt{0}, \texttt{0 1 3 2}; checked against ``$a_i$ is an ancestor of $a_j$ iff $a_i$ is the maximum between them'' on $4000$ random permutations.}
\sk{P2.04 Minimum Score of a Path Between Two Cities.}{A path may reuse roads and pass through $1$ and $n$ several times, so it can make a detour over \emph{any} road of the component of city $1$ (which contains $n$). Answer: the smallest weight of an edge in that component --- one DFS from $1$, taking the minimum over every edge scanned. \emph{Trap:} a shortest-path or ``best simple path'' search answers a different question (example~2 needs the detour $1 \to 2 \to 1$). \bigO{n + m}. Examples: $5$, $2$; checked against a search over states (city, smallest road so far) on $1886$ random graphs.}
\sk{P2.05 Count Simple Paths (ABC378 D).}{Grid graph of the empty cells. From every empty start cell, backtrack: mark the cell, try the $4$ neighbours that are empty and unmarked with one move less, un-mark on return; count when $K$ moves are done. At most $HW \cdot 4 \cdot 3^{K-1}$ calls, about $2 \cdot 10^7$ for $K = 11$; depth $K$. Samples: $2$, $0$, $218070$; checked against enumerating all sequences of $K+1$ cells on $200$ random grids up to $3 \times 3$.}
\sk{P2.06 Disappearing Permutation (CF 2086C).}{Arrows $i \to p_i$: disjoint cycles. If $p_d$ becomes $0$, position $d$ must be set to $d$; then value $d$ is no longer free for the position $j$ with $p_j = d$, which must be fixed too, and so on around the cycle --- the whole cycle of $d$ becomes the identity, while untouched cycles stay valid. Keep a \texttt{fixed} flag; each query walks the cycle of $d$ only if it is not fixed yet and adds its length. Every vertex is walked once: \bigO{n}. Samples: \texttt{1 2 3}, \texttt{2 4 4 5 5}, \texttt{4 4 4 4 7 7 7}; checked against trying all subsets of positions to reset on $1000$ random cases ($n \le 8$).}
\sk{P2.07 Prerequisites (ABC315 E).}{Arrow $i \to P_{i,j}$ (``$i$ needs $P_{i,j}$''). The books to read are exactly those reachable from $1$; a book's DFS call ends after the calls of all its prerequisites, so the \emph{post-order} (without book $1$) is a valid reading order. $N$ up to $2 \cdot 10^5$ on a chain: iterative DFS with a per-vertex iterator. \bigO{N + \sum C_i}. Samples: we print \texttt{3 5 2 4} (the sample's \texttt{5 3 4 2} is another valid order), \texttt{6 5 4 3 2}, \texttt{5}; a checker (same set as a reachability search, every prerequisite earlier) accepted all samples and $200$ random DAGs.}
\end{solutionbox}

\begin{solutionbox}[Tier 2 --- Core]
\sk{P2.08 Phase Shift (CF 1735C).}{Functional graph on letters: \texttt{nxt[y] = x} means $y$ is encrypted as $x$; the final graph must be one cycle through all $26$ letters. Scan $t$: if $x = t_i$ has no predecessor yet, try $y = \texttt{a}, \texttt{b}, \dots$ with no successor yet and $y \ne x$; reject $y$ if walking \texttt{nxt} from $x$ ends at $y$ after fewer than $26$ letters (the arrow $y \to x$ would close a short cycle). \emph{Trap:} the natural greedy ``smallest unused letter'' fails exactly there --- sample~2 must answer \texttt{ac}, not \texttt{ab}. \bigO{26^2} per new letter plus \bigO{n}. Samples: \texttt{b}, \texttt{ac}, \texttt{abcdebfadg}, and the two $26$-letter answers; checked against all circular orders of a $5$-letter alphabet (same code with $26$ replaced by $5$) on $3000$ random strings.}
\sk{P2.09 Nastya and Potions (CF 1851E).}{Arrow from a potion to each ingredient; no potion is obtainable from itself, so the graph is a DAG. $\textit{cost}(i) = 0$ for owned potions, else $\min(c_i, \sum_j \textit{cost}(e_j))$ (just $c_i$ without a recipe), computed when the DFS call of $i$ returns, memoised. Iterative post-order (chains of length $2 \cdot 10^5$); sums in \texttt{long long}. \bigO{n + \sum m_i}. Samples: \texttt{23 8 0 5 10}, \texttt{0 143 0}, \texttt{5 0 1 3 4}, \texttt{0 0 0 0}, \texttt{0 0 0 0 0 2}, \texttt{0 0 0 0 0 0}, \texttt{0 0}; checked against relaxing all recipes until nothing changes on $1500$ random DAGs.}
\sk{P2.10 Add Zeros (CF 2027C).}{The appended zeros never satisfy the condition ($a_i \ge 1$), so only the original positions matter. Position $i \ge 2$ is usable exactly when the length is $u_i = a_i + i - 1$, and it moves the length to $u_i + i - 1$: a directed edge $u_i \to u_i + i - 1$ between lengths up to $\approx 10^{12}$. Store the edges in a \texttt{map}, DFS from length $n$, answer the largest length visited. \bigO{n \log n}. Samples: $10, 11, 10, 1$; checked against a search over the actual arrays on $1500$ random cases.}
\sk{P2.11 Idiot First Search (CF 2195E).}{Let $f(v)$ be the number of moves from entering $v$ with its whole subtree blank until Bob steps to the parent: $f = 1$ for a leaf, $f(v) = f(l) + f(r) + 3$ otherwise (down left, the left tour, down right, the right tour, up). Whenever Bob climbs into an ancestor, everything there is blank again, so $\tau(k) = f(k) + f(\textit{par}(k)) + \dots + f(1)$: compute $f$ bottom-up and push the prefix sum down, all modulo $10^9+7$. \bigO{n} with an explicit stack (the tree can be a path of depth $1.5 \cdot 10^5$). Samples: \texttt{1}, \texttt{9 10 14 15 15}, \texttt{13 22 14 27 23 28 28}; checked against a literal simulation of the walk on $1200$ random trees.}
\sk{P2.12 Valera and Elections (CF 369C).}{Root at district $1$. A problem road $(\textit{par}(v), v)$ is repaired iff some elected candidate lives in the subtree of $v$. Post-order: if the road above $v$ is a problem road and nothing in $v$'s subtree is elected, elect $v$ (the deepest choice covers the most roads); propagate ``something below is elected'' to the parent. \bigO{n}. Samples: \texttt{1 / 5}, \texttt{1 / 3}, \texttt{4 / 2 3 4 5} (any order is accepted); a checker (all problem roads repaired, size equal to the minimum over all subsets) accepted $300$ random trees.}
\sk{P2.13 Subtree K-th Max (ABC239 E).}{Since $K \le 20$, each vertex returns the $20$ largest values of its subtree: merge its own value with its children's lists, sort, truncate. Answer queries from the stored lists. \bigO{20\,n \log n}; iterative post-order (a path of $10^5$ vertices). Samples: \texttt{4 5}, \texttt{9 10}, \texttt{1 10 100 1000}; checked against sorting each queried subtree on $300$ random trees with $30$ queries each.}
\sk{P2.14 Sum in the tree (CF 1098A).}{Erased sums sit at even depth. Raising $s_v$ for such a $v$ lowers the sum of all children's $a$ by one each, so take $s_v$ as large as allowed: $\min$ of the children's $s$, or $s_{\textit{par}(v)}$ for a leaf (then $a_v = 0$). Since $p_i < i$, one pass in index order fills the values, a second computes $a_v = s_v - s_{\textit{par}(v)}$; any negative $a_v$ means $-1$. \emph{Trap:} choosing $s_v = s_{\textit{par}(v)}$ for every erased vertex looks minimal but piles the cost on the children. \bigO{n}. Samples: $1, 2, -1$; checked against trying all values of the erased sums on $300$ random trees (including corrupted inputs).}
\sk{P2.15 Height of Binary Tree After Subtree Removal Queries.}{Pre-order DFS with timestamps: the subtree of $v$ is the interval $[\mathit{tin}_v, \mathit{tout}_v]$ of the pre-order, and removing it leaves exactly the vertices outside. Store the depth at each pre-order position, build prefix and suffix maxima, and answer $\max(\textit{pre}[\mathit{tin}_q], \textit{suf}[\mathit{tout}_q + 1])$. \bigO{n + m}. Examples: $[2]$, $[3, 2, 3, 2]$; checked against recomputing the height for every possible query on $3000$ random trees.}
\sk{P2.16 Game Master (CF 1608C).}{Arrow $u \to v$ if $u$ beats $v$ on some map. Player $x$ wins iff he reaches every player (eliminate them in reverse DFS order, each by the one who discovered him). Only arrows between neighbours of the two sorted orders are needed (beating is transitive along a sorted order), so the graph has $2(n-1)$ arrows. The strongest player $s$ on map~1 reaches everybody, so $x$ wins iff $x$ reaches $s$: one reverse search from $s$ (arrows from each player to the next stronger one on either map) marks all winners. \bigO{n \log n} for sorting. Samples: \texttt{0001}, \texttt{1111}, \texttt{1}; checked against a bitmask search over tournament orders on $1500$ random cases ($n \le 7$).}
\end{solutionbox}

\begin{solutionbox}[Tier 3 --- Challenge]
\sk{P2.17 Tree Pruning (CF 2018C).}{Fix the final common leaf depth $D$. A vertex $v$ can stay iff $\text{depth}(v) \le D$ and some vertex of depth $\ge D$ lies below it, i.e.\ $D \le \textit{mx}(v)$ (deepest depth in the subtree, computed bottom-up); keeping exactly these is optimal and reachable by leaf removals. Add $+1$ on $[\text{depth}(v), \textit{mx}(v)]$ with a difference array; the answer is $n - \max_D$. \bigO{n}, iterative DFS. Samples: $2, 2, 5$; checked against all root-containing subtrees on $1500$ random trees ($n \le 11$).}
\sk{P2.18 Iris and Game on the Tree (CF 2006A).}{The changes $\mathtt{10}$/$\mathtt{01}$ along a root--leaf string cancel except for the ends: the weight is non-zero iff the leaf's value differs from the root's. Count leaves $c_0, c_1$ (known) and $q$ (\texttt{?}), and $u$ = \texttt{?} vertices that are neither root nor leaf (harmless moves). Root known: $(\text{leaves of the other value}) + \lceil q/2 \rceil$. Root \texttt{?} and $c_0 \ne c_1$: Iris colours the root at once, $\max(c_0, c_1) + \lfloor q/2 \rfloor$. Root \texttt{?} and $c_0 = c_1$: nobody wants to colour the root first; the parity of $u$ decides who must, giving $c_0 + \lceil q/2 \rceil$ for odd $u$ and $c_0 + \lfloor q/2 \rfloor$ for even $u$. \bigO{n}. Samples: $2, 1, 1, 2, 1, 0$; checked against a full minimax over the game on $1200$ random trees ($n \le 8$).}
\sk{P2.19 Reachability and Tree (CF 2112D).}{Every directed edge is a good pair, so $n - 1$ pairs are free and exactly one directed path of length $2$ must exist (two would give $\ge n+1$ pairs). Pick a vertex $c$ of degree $2$ with neighbours $a, b$; orient $a \to c \to b$, make $a$ a source and $b$ a sink, and alternate sources and sinks by DFS on both sides --- then $(a, b)$ is the only extra pair. Without a degree-$2$ vertex the answer is \texttt{NO} (a length-$2$ path through a vertex of degree $\ge 3$ creates at least two extra pairs). \bigO{n}. Samples: \texttt{YES YES NO YES} with our own orientations; a checker (same edges, exactly $n$ good pairs, \texttt{NO} only if no orientation of all $2^{n-1}$ works) accepted $1200$ random trees.}
\sk{P2.20 Trapmigiano Reggiano (CF 2071C).}{Root the tree at $\textrm{en}$ and output the vertices by decreasing depth. Invariant: after the cheese at depth $d$ the mouse is at depth $\le d$ --- if it is deeper, its step towards a vertex of depth $d$ goes up, and if it is at depth $\le d$, one step changes the depth by $1$ and ends at most at $d$. After the last cheese (depth $0$) the mouse is at $\textrm{en}$. \emph{Trap:} the statement offers $-1$, but it is never needed. \bigO{n \log n}. Samples: our orders differ from the samples; a simulator of the mouse accepted the samples and $1500$ random trees ($n \le 12$, never $-1$).}
\sk{P2.21 Mad City (CF 1873H).}{$n$ vertices, $n$ edges, connected: exactly one cycle (peel degree-$1$ vertices; what remains is the cycle). On the cycle Valeriu can always run away from Marcel. Let $c$ be the cycle vertex nearest to $b$ (where he enters it); he escapes iff $\textit{dist}(b, c) < \textit{dist}(a, c)$, two BFS runs. \emph{Trap:} with equality Marcel waits at $c$; and $b$ already on the cycle gives distance $0$, so the answer is \texttt{YES} unless $a = b$. \bigO{n}. Samples: \texttt{YES NO YES NO NO YES}; checked against a retrograde solution of the full pursuit game on $1500$ random unicyclic graphs ($n \le 8$).}
\sk{P2.22 Bakry and Partitioning (CF 1592C).}{Let $X$ be the xor of all values. $X = 0$: cut any edge, both parts have equal xor. $X \ne 0$: an even number of parts would xor to $0$, so we need $3$ parts of xor $X$ (more can be merged in threes). Post-order: when the remaining xor of a subtree equals $X$, cut it off and pass $0$ upwards; we need at least two cuts and $k \ge 3$. \emph{Trap:} forgetting $k \ge 3$, or counting the root's remainder as a cut. \bigO{n}. Samples: \texttt{NO YES NO YES NO}; checked against all sets of deleted edges on $1800$ random trees ($n \le 9$).}
\sk{P2.23 Fake Plastic Trees (CF 1693B).}{Post-order greedy. Each child $c$ ends with a value $\textit{val}(c)$, and operations ending inside the children can lift $v$ by up to $S = \sum \textit{val}(c)$ for free (the arrays are non-decreasing). If $S < l_v$, start one new operation at $v$ and set $\textit{val}(v) = r_v$; else $\textit{val}(v) = \min(r_v, S)$. Since $p_i < i$, a loop from $n$ down to $1$ is the post-order. \bigO{n}. Samples: $1, 2, 2, 5$; checked against a BFS over all value vectors on $900$ random trees ($n \le 4$, $r \le 3$).}
\sk{P2.24 Uncle Bogdan and Country Happiness (CF 1388C).}{Subtree sums give $\textit{sub}_v$, the number of people passing city $v$. With $g_v$ happy ones, $h_v = 2 g_v - \textit{sub}_v$, so $g_v = (\textit{sub}_v + h_v)/2$ must be an integer in $[0, \textit{sub}_v]$; and since a mood never improves, the happy people in the children cannot exceed those in $v$: $\sum_c g_c \le g_v$. These local checks are also sufficient. Two post-order passes. \bigO{n}. Samples: \texttt{YES YES}, \texttt{NO NO}; checked against enumerating where every person's mood breaks on $1200$ random cases.}
\sk{P2.25 Edge Weight Assignment (CF 1338B).}{Root at a leaf. Minimum: if every leaf is at even depth, all weights equal work ($f = 1$); otherwise $f = 2$ is impossible and $f = 3$ always suffices (with weights $1, 2, 3$, whose xor is $0$). Maximum: start from $n - 1$ distinct weights; leaves hanging on the same vertex must have equal edge weights, so subtract $(\text{leaf neighbours} - 1)$ for every vertex with at least two of them --- everything else can be made distinct. \bigO{n}. Samples: \texttt{1 4}, \texttt{3 3}, \texttt{1 6}; checked against all assignments of weights $1..15$ on $60$ random trees ($n \le 6$).}
\sk{P2.26 Root was Built by Love, Broken by Destiny (CF 2127D).}{Every bridge crosses the river, so the graph is bipartite with the two sides as parts, and every cycle drawn this way has crossing bridges: if $m \ne n - 1$, answer $0$. In a tree, the non-leaf vertices must form a path (a caterpillar): a vertex with three non-leaf neighbours forces a crossing, giving $0$. Otherwise the spine alternates sides, the leaves of each spine vertex sit in a block between its spine neighbours in any order: $\textit{ans} = 2 \cdot (2 \text{ if the spine has} \ge 2 \text{ vertices}) \cdot \prod_v (\text{leaves at } v)!$, with $2$ for $n = 2$. \bigO{n + m}. Samples: $2, 0, 8, 12$; checked against all side assignments and orders on $1200$ random graphs ($n \le 6$).}
\end{solutionbox}

\begin{solutionbox}[Stretch --- above this chapter]
\sk{P2.27 Reachability from the Capital (CF 999E).}{Mark everything $s$ reaches. A new road may always start at $s$. For each unmarked $v$, count the unmarked vertices it reaches (one DFS each, \bigO{n(n+m)} for $n, m \le 5000$); process them by decreasing count, and whenever $v$ is still unmarked add a road $s \to v$ and mark what $v$ reaches. A vertex with the largest count is a source of the condensation, so every road added this way is necessary. Samples: $3$, $1$; checked against trying every set of new roads $s \to v$ on $300$ random graphs ($n \le 9$).}
\sk{P2.28 Tree Shuffling (CF 1363E).}{An operation at $u$ may be replaced by one at the cheapest ancestor, so push down $c_v = \min(a_v, c_{\textit{par}(v)})$. Post-order: count the wrong digits in the subtree, $x$ of type $0 \to 1$ and $y$ of type $1 \to 0$; at $v$ fix $\min(x, y)$ pairs at cost $2 c_v$ each (deeper means cheaper) and pass the rest up. Leftovers at the root mean $-1$ (different numbers of $1$s). \bigO{n}, iterative DFS. Samples: $4$, $24000$, $-1$; checked against Dijkstra over all $2^n$ digit patterns with every possible shuffle on $300$ random trees ($n \le 6$).}
\sk{P2.29 Ice Cave (CF 540C).}{Flood fill from the start over intact cells, treating the target as enterable but never walking through it. Cracked target: answer = reached. Intact target: we arrive, crack it, must step out to an intact neighbour and come back --- reached and at least two intact neighbours, counting the start cell as one if it is adjacent (we can arrive from the start directly). \emph{Trap:} start $=$ target needs just one intact neighbour; and an intact target is not enough to reach once. \bigO{nm}. Samples: \texttt{YES NO YES}; checked against a search over (position, set of cracked cells) on $300$ random grids up to $3 \times 4$.}
\sk{P2.30 Coloring Edges (CF 1217D).}{If the digraph has no cycle, one colour suffices. Otherwise two colours always suffice: colour the back edges of a DFS (edges into a grey vertex, Lesson~\ref{ssec:ch2-cycles}) with $2$ and all other edges with $1$. Tree, forward and cross edges go from a vertex with a larger exit time to one with a smaller exit time, back edges the other way, so neither class can close a cycle. \emph{Trap:} trying to minimise the number of back edges is unnecessary --- only $k$ is minimised. \bigO{n + m}. Samples: \texttt{1 / 1 1 1 1 1}, \texttt{2 / 1 1 2}; a checker ($k$ equal to $1$ iff the graph is acyclic, every colour class acyclic) accepted $300$ random digraphs.}
\end{solutionbox}
\endgroup

%------------------------------------------------------------
\subsection{Answers to the Self-Check Quiz}
\label{ssec:ch2-quiz-answers}

\begin{solutionbox}[Answers]
\begin{enumerate}[leftmargin=*]
    \item The path from the start vertex to the current vertex, along the tree edges through which DFS went deeper --- exactly the gray vertices (Lesson~\ref{ssec:ch2-call-stack}).
    \item It never knows when a vertex is finished: a vertex is popped once and all its neighbours are pushed at once, so there is no moment ``all of $u$'s neighbours are done''. Its stack is also not a path and may hold a vertex several times.
    \item Two DFS intervals $[\tin{}, \tout{}]$ are either disjoint or nested, never partially overlapping; nested means ancestor and descendant. $\tin{u} < \tin{v} < \tout{v} < \tout{u}$: $v$ is a descendant of $u$.
    \item Tree (target white), back (target gray), forward (target black, $\tin{u} < \tin{v}$), cross (target black, $\tin{v} < \tin{u}$).
    \item If $\{u, v\}$ is an edge and $u$ is entered first, $u$ scans $v$ while gray: either $v$ is still white (tree edge) or it was entered meanwhile and is a descendant, so the edge is first seen from $v$'s side as a back edge to the gray $u$.
    \item A black target was finished earlier; it only means two paths lead there. Example: $1 \to 2 \to 3$ and $1 \to 3$ --- vertex 3 is black when $1 \to 3$ is examined, and there is no cycle.
    \item Two parallel edges between $p$ and $v$ form a cycle of length 2. Skipping only the edge used to enter $v$ lets the second one be found as a back edge; skipping the vertex $p$ hides it.
    \item A forest has $n - c$ edges, so a cycle exists iff $m > n - c$.
    \item The children's sizes $\mathrm{sz}[c]$ for every child $c$ of $v$, and the part above, $n - \mathrm{sz}[v]$ (if $v$ is not the root).
    \item The block $\mathrm{pos}[u], \mathrm{pos}[u] + 1, \dots, \mathrm{pos}[u] + \mathrm{sz}[u] - 1$.
    \item The mark of a vertex is removed when it is left. Then the marks mean ``on the current path'' and the search tries every simple path --- exponentially many in general, instead of \VE.
    \item Reverse every edge and run one search from $t$ in $G^R$; the marked vertices are exactly those that can reach $t$ in $G$.
\end{enumerate}
\end{solutionbox}

